\subsection{Metric Landmark Extraction and Adaptive Smoothing}\label{subsec:landmark_smoothing}

For each detected hand, MediaPipe returns two distinct landmark sets, and
\texttt{hand\_kinematics\_3D} uses both for complementary purposes. The
\emph{world} landmarks are expressed in metric units with the wrist at the
origin, and so describe the hand's intrinsic shape and proportions; the
\emph{image} landmarks are normalised to the frame and so describe where the hand
sits on screen. The \texttt{Hand3D} class converts these into the input the rest
of the pipeline expects.

The hand shape is taken from the world landmarks, scaled by a constant factor, and
smoothed with the 1\euro{} filter~(\autoref{subsec:one_euro_filter}). A separate filter
instance is maintained for each hand so that the two are smoothed independently.

Because the world landmarks are wrist-relative, they carry no information about how
far the hand is from the camera. That information is restored separately, in an
\emph{anchor} that locates the hand in the scene. The anchor combines the wrist's
image-space position with a depth estimate produced by \texttt{\_estimate\_scale},
which compares the mean image-space length of a stable subset of bones against
their mean metric length in world space. The subset is fixed in advance --- the
thumb chain, the index and pinky bases, and the row of knuckles --- and the ratio
of the two means serves as a monocular depth cue, since a hand closer to the camera
projects to larger image-space distances. The anchor is passed through its own
1\euro{} filter, again one instance per hand. \texttt{get\_3d\_coordinates}
returns the filtered shape and the filtered anchor as two separate sets, so that a
metrically correct hand can later be placed at the estimated location.
